% Pursuing Stacks (cotes 134-2 to 134-8): archive cote 134-4, site page 67 (PDF page 68), typescript page 155.
% typescript. Draft compared: ps3.tex:3437-3520.
% AI-transcribed by Claude Opus 5.5 (claude-opus-5-5, adaptive thinking, effort high), Message Batches API, batch ps2 (msgbatch_01NVwr1Yuo8qJf5rykeGBXPK), 2026-11-12; not yet checked by eye.
\scanpage{134-4/67}
\mspage{155}

f) Let \hand{u}: \unsure{X'}$\to$ \unsure{X} be a map in(Cat)such that for any a in X, the
induced category $X'_{/a}$ be $\underline{W}$-aspheric i.e. $X'_{/a}\to e$ is in $\underline{W}$
(or what amounts to the same if we assume d) e.g. if we assume
the stronger condition c)\deleted{\illegible{1}}\hand{,} the induced map $X'_{/a}\to X_{/a}$ is in $\underline{W}$).
Then \hand{u} is in $\underline{W}$ .

If \hand{u}\inserted{: $X'\to X$} satifies the condition stated above, namely that after any base
change $X_{/a}\to X$ , the corresponding map $\hand{u}_{/a}$ is in $\underline{W}$ , we'll say that \hand{u}
is \uline{weakly} $\underline{W}$-aspheric (whereas \xout{\unsure{"asph}} $\underline{W}$-aspheric" means that after \uline{any}
base change $Y_{\unsure{\hand{x}}}\to X$, the corresponding \xout{\illegible{1}}\inserted{\unsure{I}} $f_Y$ is in $\underline{W}$ ). Thus con-%
dition f) can be stated, saying that \uline{a weakly} $\underline{W}$-\uline{aspheric} \uline{map} \uline{in} (Cat)
\uline{is in} $\underline{W}$ .

For \xout{\unsure{usin}} making use of this latter assumption on $\underline{W}$ , we have to
look at how the induced categories for the functor (12) look like, which
functor (I recall) induced a bijection on objects. \deleted{which}\inserted{These} can be describ\hand{\unsure{ed}}
as pairs (a,p) , with a in A and p a map in (Cat)
\[ p:\ i(a)\to C \quad . \]
An easy computation shows the

\uline{Lemma} \over{Det}{Let} (a,p) as above. THe induced category $({A}_{/C})_{/(a,p)}$ (for
the functor (12)) is canonically isomorphic to the induced category
$A_{/G}$ , where G is the fibered product in A\^{} displayed in the diagram
\diagram{
(13)\qquad
\begin{tikzcd}
G \arrow[r, "\xout{\unsure{x}}"] \arrow[d] & a \arrow[d, "\hand{\unsure{p}}"] \\
i^{*}(\underline{Fl}(C)) \arrow[r] & i^{*}(C)
\end{tikzcd}\quad ,
}
where
\[ \underline{Fl}\,(C) \overset{\text{def}}{=} \underline{Hom}\,(\hand{\Delta}_1,C\ ) \]
and where the second horizontal arrow in (13) is the $i^{*}$ -transform of the
target map in (Cat)
\[ \underline{Fl}(C) \xrightarrow{\ t\ } C \qquad . \]

\uline{Corollary} In order for (12) to be weakly $\underline{W}$-aspheric, it is ns that
the map
\[ \hand{(14)} \qquad i^{*}(t)\ :\ i^{*}(\underline{Fl}(C))\longrightarrow i^{*}(C) \]
in A\^{} be $\underline{W}_A$ -aspheric (i.e. be "universally in $\underline{W}_A$").

To make the meaning of the latter condition clear, it should be
noted that the condition f) on $\underline{W}$ gu\unsure{arantees} precisely that for a
map\inserted{u:} $F'\to F$ in A\^{} (\over{for}{A} any category \deleted{\illegible{1}} ) to be \xout{in} $W_A$-aspheric, it is
ns that the corresponding map $i_A(u)$ in (Cat) be weakly $\underline{W}$-aspheric -
the kind of thing we have been \margin{transposition marks around ``using'' and ``constantly'': read ``constantly using''}using constantly before of course, when

\margin{67}
